% ============================================================
% CAPITOLUL 7 — APLICAȚIA WEB E-MEDICINE
% ============================================================

\chapter{Aplicația Web E-Medicine}
\label{ch:aplicatie}

\section{Descrierea funcționalităților aplicației}
\label{sec:functionalitati}

\todosec{Prezentare generală a aplicației: ecranele principale, fluxurile
de utilizare cheie. Screenshot-uri ale interfeței finale.}

\section{Backend Go — Implementare detaliată}
\label{sec:backend_go}

\subsection{Structura proiectului Go}
\label{subsec:go_structura_proiect}

\todosec{Arborele de directoare al proiectului Go, convenții de denumire,
separarea responsabilităților pe pachete: api, db, runner, models, config.}

\subsection{Implementarea endpointurilor REST}
\label{subsec:go_endpoints}

\todosec{Codul Go pentru endpointurile principale, utilizând framework-ul ales
(Gin/Echo). Demonstrarea pattern-ului: handler -> service -> repository -> DB.}

\subsection{Gestionarea execuției algoritmilor}
\label{subsec:go_runnere}

\todosec{Implementarea runner-elor asincrone: pool de goroutine, channel-uri
pentru rezultate, context cancellation pentru timeout-uri.
Codul Go pentru lansarea subprocess-urilor R, Julia, Python, SAS.}

\subsection{Conexiunea la YugabyteDB}
\label{subsec:go_db}

\todosec{Driverul \texttt{pgx} pentru PostgreSQL/YSQL, connection pooling,
query-uri parametrizate pentru prevenirea SQL injection.}

\subsection{Autentificare și securitate}
\label{subsec:go_securitate}

\todosec{JWT middleware, CORS configuration, rate limiting, validarea
input-ului, HTTPS cu certificate TLS.}

\section{Frontend Vue.js — Implementare detaliată}
\label{sec:frontend_vue}

\subsection{Structura componentelor și pagini}
\label{subsec:vue_componente}

\todosec{Arborele de componente, structura directorului \texttt{src/},
tipare de compunere (composables reutilizabile).}

\subsection{Dashboard epidemiologic}
\label{subsec:vue_dashboard}

\todosec{Componenta Dashboard: metrici cheie (R0 curent, vârf de infectați,
dimensiunea finală estimată), grafice de sumar, ultimele rulări.}

\subsection{Vizualizarea graficelor SIR}
\label{subsec:vue_grafice_sir}

\todosec{Componenta SIRChart cu ECharts/Chart.js: afișarea traiectoriilor
S(t), I(t), R(t), zoom interactiv, export PNG/SVG. Codul Vue 3 (Composition API).}

\subsection{Comparatorul de rezultate între limbaje}
\label{subsec:vue_comparator}

\todosec{Pagina de comparare: suprapunerea graficelor din rulări multiple
(R, Julia, Python, SAS pe același set de date), tabelul de diferențe numerice,
graficul timpilor de execuție.}

\subsection{Interfața pentru seturile de date}
\label{subsec:vue_datasets}

\todosec{CRUD pentru seturile de date: upload CSV, previzualizare, editarea
metadatelor, ștergere. Validarea client-side a formularelor.}

\section{Fluxuri de utilizare (use cases)}
\label{sec:use_cases}

\todosec{Descrierea pas-cu-pas și screenshot-uri pentru:
(1) Încărcarea unui set de date nou,
(2) Declanșarea și monitorizarea unei simulări,
(3) Compararea rezultatelor între limbaje,
(4) Exportul rezultatelor.}

\section{Testare și validare funcțională}
\label{sec:testare_functionala}

\todosec{Unit teste Go (testing package), integration tests (testcontainers),
end-to-end tests (Playwright/Cypress pentru Vue.js). Acoperire de cod.}
